% ============================================================
%  Section 3 — Methods
%  "Transfer Learning from Lithium-Ion to Zinc-Ion Batteries:
%   Cross-Chemistry Capacity Prediction with Limited Target Data"
% ============================================================

\section{Methods}

\subsection{Data preprocessing and quality control}
\label{sec:preprocessing}

Raw cycling data were drawn from three publicly available sources: (i) the
Severson/MATR dataset comprising 83 lithium iron phosphate (LFP) cells
\cite{severson2019}, (ii) 12 lithium cobalt oxide (LCO) cells from the CALCE
repository, and (iii) 130 zinc-ion manganese dioxide (Zn/MnO\textsubscript{2})
cells from the BatteryLife dataset \cite{tan2024}, which spans three
experimental batches collected under distinct temperature conditions.
Together, the 95 Li-ion cells form the \emph{source} task ($t=0$) and the
130 Zn-ion cells form the \emph{target} task ($t=1$).

All cells were passed through a ten-step quality-control pipeline prior to
any feature extraction.
(1) Cells with fewer than 50 recorded cycles were excluded, as the early-cycle
window used for feature extraction (cycles 2--100) would otherwise be
incomplete.
(2) Individual cycles whose measured discharge capacity fell below 95\% of the
cell's nominal rated capacity were flagged as incomplete---typically arising
from mid-test interruptions or instrumentation dropouts---and removed from
further analysis.
(3) Voltage and capacity traces were checked for monotonicity within each
half-cycle and non-physical discontinuities were corrected by linear
interpolation between flanking valid measurements.
(4) Duplicate cycle indices, which occasionally appear in raw BatteryLife
exports, were deduplicated by retaining the record with the higher measured
capacity.
(5) Cells missing full voltage--capacity ($V$--$Q$) curves for at least 10
cycles in the early window were excluded from the FPCA pathway and assigned
only the scalar $\Delta Q_{\mathrm{var}}$ feature, with FPCA scores imputed
to the task-level mean; such cells constituted fewer than 3\% of the corpus.
(6) Cycle-life labels were defined as the first cycle at which normalised
discharge capacity fell and remained below the end-of-life (EOL) threshold;
EOL was set to 80\% of the reference capacity $Q_{\mathrm{ref}}$ (defined in
Section~\ref{sec:activation}).
(7) Cells that never reached the EOL threshold within the available cycling
record were treated as right-censored and excluded from supervised regression
training, though they were retained for qualitative analysis.
(8) Current-rate (C-rate) metadata were verified; the small number of CALCE
cells cycled at non-standard rates were normalised to equivalent cycle
counts via Ah-throughput scaling before merging with the Severson corpus.
(9) All capacity values were expressed as a fraction of $Q_{\mathrm{ref}}$
(see Section~\ref{sec:activation}) for cross-cell comparability.
(10) The final dataset was split into a fixed held-out test set of 10 Zn-ion
cells, drawn to be representative of all three BatteryLife batches, with the
remainder used for training and validation sweeps (Section~\ref{sec:protocol}).

A critical design decision was that the functional principal component
analysis (FPCA) basis was \emph{fitted exclusively on Li-ion data} and
Zn-ion curves were subsequently \emph{projected} onto this fixed basis
without refitting (Section~\ref{sec:fpca}).  This asymmetry is intentional:
refitting the basis on Zn-ion data would destroy the shared representational
space through which knowledge is transferred, because any linear relationship
between the source and target latent features would be lost after independent
rotations.

\subsection{Activation period detection}
\label{sec:activation}

Unlike Li-ion cells, Zn/MnO\textsubscript{2} batteries frequently exhibit a
monotonic \emph{capacity increase} during an initial activation period,
attributed to progressive electrolyte wetting and electrochemical conditioning
of the MnO\textsubscript{2} cathode \cite{tan2024}.  This behaviour
invalidates the standard convention of normalising capacity by the
first-cycle discharge value: applying that convention would assign
normalised capacities greater than unity to mid-life cycles, introducing
data leakage because the true capacity maximum $Q_{\max}$ is observed
\emph{during} the experiment, not before it.  Treating $Q_{\max}$ as
the reference would likewise leak future information into the training
features.

To determine the cycle $N_{\mathrm{act}}$ at which sustained monotonic
degradation begins, we applied the Pruned Exact Linear Time (PELT)
changepoint detection algorithm \cite{truong2020} from the
\texttt{ruptures} library, using an $\ell_2$ (least-squares) cost model and
a penalty term calibrated by the Bayesian information criterion.
A candidate changepoint was accepted only if the post-break segment spanned
at least 10 consecutive cycles of non-increasing capacity, ensuring that
brief transient dips did not incorrectly terminate the activation window.
When PELT detected no statistically significant changepoint---occurring in
approximately 8\% of Zn-ion cells that showed monotonic degradation from
cycle 1---$N_{\mathrm{act}}$ was set to $\operatorname{argmax}_n Q(n)$ as a
deterministic fallback.  For all Li-ion cells, no activation period was
observed and $N_{\mathrm{act}}$ was fixed at cycle 1.

The reference capacity was defined as $Q_{\mathrm{ref}} = Q(N_{\mathrm{act}})$
and normalised capacity as
\begin{equation}
  Q_{\mathrm{norm}}(n) = \frac{Q\!\left(N_{\mathrm{act}} + n\right)}{Q_{\mathrm{ref}}},
  \qquad n = 1, 2, \ldots
  \label{eq:qnorm}
\end{equation}
Cycle-life labels were then computed as the number of cycles $n$ elapsed from
$N_{\mathrm{act}}$ until $Q_{\mathrm{norm}}(n)$ first dropped below 0.80 and
remained below that threshold, yielding values in the range 50--1662 for
Zn-ion cells and 88--2237 for Li-ion cells.  All models were trained on
log-transformed cycle-life values to stabilise variance across the wide
dynamic range.

\subsection{Functional principal component analysis}
\label{sec:fpca}

Differential capacity ($\mathrm{d}Q/\mathrm{d}V$) curves encode
electrochemical phase-transition signatures whose shifts with aging have
demonstrated predictive value for Li-ion degradation \cite{severson2019}.
We adopted these curves as the primary spectral feature carrier.

For Severson LFP cells, $Q(V)$ data from cycles 2--100 were linearly
interpolated onto a uniform 1000-point voltage grid spanning 2.0--3.5\,V,
and $\mathrm{d}Q/\mathrm{d}V$ was computed by central finite differences.
For Zn-ion cells, raw ($V$, $Q$) point arrays were interpolated onto a
500-point grid spanning 0.8--1.8\,V prior to differentiation, reflecting
the narrower electrochemical window of the Zn/MnO\textsubscript{2} couple.
In both cases, Savitzky--Golay smoothing (window length 11, polynomial
order 3) was applied to suppress numerical noise before differentiation.

FPCA was performed using the \texttt{scikit-fda} library
\cite{ramos2024} on an \texttt{FDataGrid} constructed from \emph{Li-ion
curves only}.  The $K = 3$ retained functional principal components explained
96.1\% of variance in the Li-ion corpus (PC1: 83.0\%, PC2: 10.5\%,
PC3: 2.6\%), and the component loadings were held fixed thereafter.

To project Zn-ion curves onto the Li-ion basis, the Zn-ion voltage grid was
first mapped to the normalised domain $[0, 1]$ by a linear rescaling
$v \mapsto (v - v_{\min})/(v_{\max} - v_{\min})$; the Li-ion grid was
rescaled identically.  This common normalised domain permitted direct
evaluation of the Li-ion eigenfunctions at Zn-ion grid points via cubic
spline interpolation, after which the standard
\texttt{FPCA.transform()} call was applied to obtain three projection
scores.  This procedure explicitly assumes that, under a common normalised
voltage axis, the degradation modes of the two chemistries share
\emph{sufficient} geometric overlap for the Li-ion loadings to capture
informative variance in the Zn-ion curves---an assumption validated empirically
by the positive transfer results in Section~4.

The final feature vector for each cell was
\begin{equation}
  \mathbf{x} = \bigl[\,\phi_1,\;\phi_2,\;\phi_3,\;\Delta Q_{\mathrm{var}}\bigr]^{\!\top} \in \mathbb{R}^4,
  \label{eq:features}
\end{equation}
where $\phi_k$ are the FPCA projection scores and
$\Delta Q_{\mathrm{var}} = \mathrm{Var}\!\left(Q_{\mathrm{norm}}(2) - Q_{\mathrm{norm}}(100)\right)$
is the variance of the capacity difference curve, a scalar feature previously
shown to be predictive of long-term degradation \cite{severson2019}.
Coulombic efficiency (CE) was \emph{deliberately excluded} from the feature
set.  In Li-ion cells, CE primarily reflects solid-electrolyte interphase
(SEI) growth, a continuous parasitic film formation mechanism; in
Zn-ion cells, CE instead reflects the competing effects of hydrogen evolution
reaction (HER) and dead-zinc accumulation, which are mechanistically distinct
and exhibit different cycle-number dependence.  An ablation study
(Section~4.4) confirmed that including CE produced negative transfer,
increasing RMSE on Zn-ion test cells relative to a Zn-ion-only baseline.

\subsection{Multi-task Gaussian process with intrinsic coregionalization}
\label{sec:mtgp}

Let $\mathcal{D}_s = \{(\mathbf{x}_i^{(s)}, y_i^{(s)})\}_{i=1}^{N_s}$ and
$\mathcal{D}_t = \{(\mathbf{x}_j^{(t)}, y_j^{(t)})\}_{j=1}^{N_t}$ denote the
source (Li-ion) and target (Zn-ion) training sets, where $y$ denotes
log-cycle-life.  We augmented each input with a discrete task index
$t \in \{0, 1\}$ and placed a Gaussian process (GP) prior jointly over both
tasks:
\begin{equation}
  f(\mathbf{x}, t) \sim \mathcal{GP}\!\left(0,\; k\!\left((\mathbf{x},t),\,(\mathbf{x}',t')\right)\right).
  \label{eq:gp_prior}
\end{equation}

The covariance kernel was structured under the \emph{intrinsic
coregionalization model} (ICM):
\begin{equation}
  k\!\left((\mathbf{x},t),\,(\mathbf{x}',t')\right)
  = B_{t,t'} \cdot k_{\mathbf{x}}(\mathbf{x},\mathbf{x}'),
  \label{eq:icm}
\end{equation}
where $\mathbf{B} \in \mathbb{R}^{2 \times 2}$ is a learned positive
semi-definite task covariance matrix and $k_{\mathbf{x}}$ is a shared
input-space kernel.  The ICM factorisation assumes that cross-task
correlations are \emph{stationary} in input space---i.e., the relative
similarity between two observations from different tasks depends only on
their feature-space distance, not on their absolute location---which is
appropriate here given the common normalised feature representation of
Eq.~\eqref{eq:features}.

The input-space kernel was chosen as Mat\'{e}rn-$\tfrac{5}{2}$:
\begin{equation}
  k_{\mathbf{x}}(\mathbf{x},\mathbf{x}') =
  \left(1 + \frac{\sqrt{5}\,r}{\ell} + \frac{5r^2}{3\ell^2}\right)
  \exp\!\left(-\frac{\sqrt{5}\,r}{\ell}\right),
  \label{eq:matern52}
\end{equation}
where $r = \|\mathbf{x} - \mathbf{x}'\|_2$ and $\ell > 0$ is the
learnable length-scale.  The Mat\'{e}rn-$\tfrac{5}{2}$ kernel was preferred
over the squared-exponential (RBF) kernel because it encodes a weaker
smoothness assumption (twice differentiable), which is physically more
appropriate for battery degradation trajectories that can exhibit
sharp transitions near EOL.

The full set of trainable hyperparameters was
$\boldsymbol{\theta} = \{\mathbf{B}, \ell, \sigma^2_{\mathrm{noise}}\}$,
where $\sigma^2_{\mathrm{noise}}$ is an isotropic observation noise variance
shared across tasks.  The matrix $\mathbf{B}$ was parameterised as
$\mathbf{B} = \mathbf{L}\mathbf{L}^{\top}$ with $\mathbf{L}$ lower triangular,
ensuring positive semi-definiteness throughout optimisation.
Hyperparameters were jointly optimised by maximising the log marginal likelihood
\begin{equation}
  \log p(\mathbf{y} \mid \mathbf{X}, \boldsymbol{\theta}) =
  -\tfrac{1}{2}\mathbf{y}^{\top}\mathbf{K}_{\boldsymbol{\theta}}^{-1}\mathbf{y}
  -\tfrac{1}{2}\log|\mathbf{K}_{\boldsymbol{\theta}}|
  -\tfrac{N}{2}\log 2\pi,
  \label{eq:mll}
\end{equation}
where $\mathbf{K}_{\boldsymbol{\theta}}$ is the $(N_s + N_t) \times (N_s + N_t)$
joint covariance matrix evaluated at all training inputs and task indices,
and $N = N_s + N_t$.  Optimisation used the Adam algorithm with a learning
rate of $10^{-2}$ for 500 steps; we verified convergence by confirming that
the change in marginal likelihood across the final 50 steps was less than
$10^{-4}$ in all Monte Carlo trials.  The implementation relied on
\texttt{GPyTorch} \cite{gardner2018}, using \texttt{IndexKernel} for
$\mathbf{B}$ combined with \texttt{MaternKernel} via a \texttt{ProductKernel}.

At prediction time, the posterior mean and variance for a new Zn-ion cell
with features $\mathbf{x}_*$ (task index $t=1$) are given by the standard GP
posterior:
\begin{align}
  \mu_* &= \mathbf{k}_*^{\top}\mathbf{K}_{\boldsymbol{\theta}}^{-1}\mathbf{y},
  \label{eq:posterior_mean} \\
  \sigma^2_* &= k(\mathbf{x}_*,\mathbf{x}_*) - \mathbf{k}_*^{\top}\mathbf{K}_{\boldsymbol{\theta}}^{-1}\mathbf{k}_*,
  \label{eq:posterior_var}
\end{align}
where $\mathbf{k}_* = [k((\mathbf{x}_*,1),(\mathbf{x}_i,t_i))]_i$ is the
vector of covariances between the test point and all training points.  The
posterior standard deviation $\sigma_*$ thus provides calibrated predictive
uncertainty that reflects both data scarcity (small $N_t$) and structural
uncertainty in the inter-task relationship.

We chose a GP formulation over neural transfer approaches (e.g., fine-tuning
a multi-layer perceptron trained on Li-ion data) for three reasons specific
to the low-data regime.  First, at $N_t \in \{2,5,10\}$ labeled Zn-ion
cells, GPs provide well-calibrated uncertainty estimates without requiring
held-out validation data for early stopping.  Second, the ``freeze early
layers'' paradigm of neural transfer learning is not applicable to
shallow tabular models where all parameters are jointly tuned; GP
marginal likelihood optimisation achieves analogous regularisation
automatically by integrating over function-space complexity.  Third, the
ICM's off-diagonal element $B_{01}$ is a directly interpretable quantity:
its sign and magnitude encode whether source knowledge transfers
constructively or destructively to the target, and can be inspected
to diagnose negative transfer.

\subsection{Experimental protocol}
\label{sec:protocol}

To characterise how prediction accuracy degrades as labeled Zn-ion training
data become scarcer, we conducted a systematic sweep over target set sizes
$N_t \in \{2, 5, 10, 20, 40\}$.  For each value of $N_t$, we drew 200
independent Monte Carlo random subsamples (without replacement) from the
pool of non-held-out Zn-ion cells.  The sampled $N_t$ cells formed the
target training set; the remaining non-held-out cells served as an
in-distribution validation set for monitoring generalisation.  All 95 Li-ion
source cells were included in every trial.  Performance was evaluated on the
fixed 10-cell Zn-ion test set, which was never used during hyperparameter
optimisation or model selection in any trial.

The primary metric was root mean squared error (RMSE) in cycle units,
obtained by computing residuals after inverting the log transform applied to
cycle-life labels.  For each $N_t$, we report the median RMSE and
interquartile range (IQR) across the 200 Monte Carlo trials to provide
a robust picture of performance variability under different random draws of
the scarce target training data.

Statistical significance of performance differences between the MT-GP and each
baseline (Section~\ref{sec:baselines}) was assessed using the Wilcoxon
signed-rank test \cite{wilcoxon1945}, pairing trials by random seed so that
both models saw identical training sets.  $p$-values were corrected for
multiple comparisons across $N_t$ values and baseline models using the
Bonferroni method; results are reported as significant at the corrected
$\alpha = 0.05$ level.

\subsection{Baselines}
\label{sec:baselines}

Three baselines were evaluated against the MT-GP.

\textbf{GP-Direct.} A single-task GP trained solely on the $N_t$ available
Zn-ion cells, using the same Mat\'{e}rn-$\tfrac{5}{2}$ kernel and identical
hyperparameter optimisation procedure.  This baseline quantifies the
improvement attributable to the Li-ion source data specifically, holding
the model class and feature representation constant.

\textbf{Ridge regression.} An $\ell_2$-regularised linear regressor fit on
the Zn-ion training cells, with regularisation strength $\alpha$ selected
by leave-one-out cross-validation over the grid
$\alpha \in \{10^{-4}, 10^{-3}, \ldots, 10^{4}\}$.  Ridge regression serves
as a linear lower-bound representative of classical, non-transfer approaches
and contextualises whether GP non-linearity provides value at very small
$N_t$.

\textbf{Shuffled-label control.} The MT-GP trained with Li-ion cycle-life
labels randomly permuted prior to each trial, while all other aspects of
the pipeline (features, kernel, optimisation) remain unchanged.  This
control tests the specific hypothesis that MT-GP performance improvements
arise from \emph{informative} source supervision, rather than from the
mere regularisation effect of additional data points with arbitrary labels.
A failure of MT-GP to outperform the shuffled control would indicate that
the Li-ion source knowledge contributes nothing beyond increasing the
effective training set size; a clear advantage over the shuffled control is
necessary (though not sufficient) evidence of genuine cross-chemistry
knowledge transfer.
